\textbf{Solutions -- Part 1}

\begin{description}
\item[(a)] If we plot given coordinates in XY, XZ and YZ planes we'll see
  the following pictures. See Figures 2, 3 and 4 (0.5p for each graph).
  % FIGURE NEEDED: XY, XZ and YZ scatter plots of the 16 satellite
  % coordinates ("viewpoint of a cluster from observational site"),
  % 2022_da_2_sol.pdf pages 5-6
  By observing Figure 2 we can deduce that the group is edge-on and its
  normal lies in XY plane. $X'$ axis is along line of sight so it can be
  calculated by normalizing
  $8.3\hat{i} + 6.2\hat{j} + 1.1\hat{k}$. So
  \[ \hat{x}' = 0.79667866\,\hat{i} + 0.59510936\,\hat{j}
     + 0.10558392\,\hat{k} \]
  (0.5p) $Z'$ axis lays in XY plane and is perpendicular to $X'$ so it
  can be calculated by normalizing $\hat{x}'\times\hat{k}$
  \[ \hat{z}' = 0.59845448\,\hat{i} - 0.80115681\,\hat{j} \]
  \hfill(2p)

\item[(b)] One can choose $x'$ as line of sight (found in part (a)) (1p):
  \[ \hat{x}' = 0.79667866\,\hat{i} + 0.59510936\,\hat{j}
     + 0.10558392\,\hat{k} \]
  $z'$ can be chosen as halo disk normal (found in part (a)) (1p):
  \[ \hat{z}' = 0.59845448\,\hat{i} - 0.80115681\,\hat{j} \]
  And finally for $y'$ (1p):
  \[ \hat{y}' = \hat{x}'\times\hat{z}'
     = 0.08458928\,\hat{i} + 0.06318717\,\hat{j} - 0.9944104\,\hat{k} \]

\item[(c)] Each coordinate in $X'Y'Z'$ system can be found by following
  equations
  \begin{align*}
  x' &= \hat{x}'\cdot(\vec{r} - \vec{r}_c)\\
  y' &= \hat{y}'\cdot(\vec{r} - \vec{r}_c)\\
  z' &= \hat{z}'\cdot(\vec{r} - \vec{r}_c)
  \end{align*}
  Where $\hat{x}'$, $\hat{y}'$ and $\hat{z}'$ are unit vectors of local
  $X'Y'Z'$ system, $\vec{r}$ is coordinate of a galaxy and $\vec{r}_c$
  is coordinate of the center.

  This values are calculated in the following table (0.5p for each
  correct $x'$, $y'$, $z'$ tuple):

  \begin{center}
  \begin{tabular}{rrrrrrr}
  & $x - x_c$ & $y - y_c$ & $z - z_c$ & $x'$ & $y'$ & $z'$\\
  1 & 0.101124 & 0.109271 & 0.294476 & 0.176684 & $-0.277372$ & $-0.027025$\\
  2 & 1.582909 & 0.989392 & 0.405895 & 1.892722 & $-0.207212$ & 0.154641\\
  3 & 0.582735 & 0.213285 & 1.126426 & 0.710112 & $-1.057359$ & 0.177866\\
  4 & 0.144328 & 0.369381 & 1.338567 & 0.476137 & $-1.295537$ & $-0.209558$\\
  5 & $-0.517770$ & $-0.151687$ & $-1.458202$ & $-0.656729$ & 1.396668 & $-0.188336$\\
  6 & 0.480059 & 0.105110 & 1.362674 & 0.588881 & $-1.307808$ & 0.203084\\
  7 & $-0.309855$ & $-0.319116$ & $-0.568033$ & $-0.496739$ & 0.518483 & 0.070227\\
  8 & 0.239541 & 0.188218 & $-0.730736$ & 0.225694 & 0.758807 & $-0.007438$\\
  9 & $-0.324800$ & $-0.169551$ & 0.116378 & $-0.347375$ & $-0.153915$ & $-0.058541$\\
  10 & $-1.227659$ & $-0.768684$ & $-0.612379$ & $-1.500158$ & 0.456538 & $-0.118862$\\
  11 & $-0.263209$ & $-0.234869$ & $-0.485227$ & $-0.400698$ & 0.445410 & 0.030649\\
  12 & 0.381200 & 0.282515 & $-0.365581$ & 0.433222 & 0.413634 & 0.001792\\
  13 & $-1.185021$ & $-0.941303$ & 0.302593 & $-1.472310$ & $-0.460620$ & 0.044950\\
  14 & 1.286640 & 0.930147 & $-0.777889$ & 1.496445 & 0.941150 & 0.024802\\
  15 & $-0.106760$ & $-0.095804$ & $-0.184997$ & $-0.161600$ & 0.168879 & 0.012863\\
  16 & $-0.686865$ & $-0.352750$ & 0.566244 & $-0.697350$ & $-0.643470$ & $-0.128449$
  \end{tabular}
  \end{center}

\item[(d)] Spherical coordinates can by calculated by following
  relations:
  \[ r = \sqrt{x'^2 + y'^2 + z'^2} \]
  \[ \phi = \arctan\frac{y'}{x'}\ \text{and signs of $x'$ and $y'$} \]
  \[ \theta = \arccos\frac{z}{r} \]
  Calculated values are given in the following table (0.5p for each
  correct $r$, $\phi$, $\theta$ tuple):

  \begin{center}
  \begin{tabular}{rrrrrrr}
  & $x'$ & $y'$ & $z'$ & $r$ & $\phi$ & $\theta$\\
  1 & 0.176684 & $-0.277372$ & $-0.027025$ & 0.329974 & 5.279566 & 1.652789\\
  2 & 1.892722 & $-0.207212$ & 0.154641 & 1.910301 & 6.174141 & 1.489757\\
  3 & 0.710112 & $-1.057359$ & 0.177866 & 1.286042 & 5.303793 & 1.432047\\
  4 & 0.476137 & $-1.295537$ & $-0.209558$ & 1.396079 & 5.064586 & 1.721471\\
  5 & $-0.656729$ & 1.396668 & $-0.188336$ & 1.554814 & 2.010330 & 1.692226\\
  6 & 0.588881 & $-1.307808$ & 0.203084 & 1.448581 & 5.135477 & 1.430138\\
  7 & $-0.496739$ & 0.518483 & 0.070227 & 0.721461 & 2.334780 & 1.473302\\
  8 & 0.225694 & 0.758807 & $-0.007438$ & 0.791695 & 1.281697 & 1.580191\\
  9 & $-0.347375$ & $-0.153915$ & $-0.058541$ & 0.384430 & 3.558678 & 1.723672\\
  10 & $-1.500158$ & 0.456538 & $-0.118862$ & 1.572587 & 2.846171 & 1.646452\\
  11 & $-0.400698$ & 0.445410 & 0.030649 & 0.599906 & 2.303400 & 1.519685\\
  12 & 0.433222 & 0.413634 & 0.001792 & 0.598981 & 0.762273 & 1.567805\\
  13 & $-1.472310$ & $-0.460620$ & 0.044950 & 1.543336 & 3.444801 & 1.541667\\
  14 & 1.496445 & 0.941150 & 0.024802 & 1.767973 & 0.561416 & 1.556768\\
  15 & $-0.161600$ & 0.168879 & 0.012863 & 0.234094 & 2.334171 & 1.515821\\
  16 & $-0.697350$ & $-0.643470$ & $-0.128449$ & 0.957522 & 3.886829 & 1.705349
  \end{tabular}
  \end{center}

\item[(e)] Using radii calculated above and number of galaxies for each
  data one can calculate number of galaxies in each bin. Range of bin 1
  is (0.2, 0.425), bin 2 -- (0.425, 0.65) and so on bin 8 --
  (1.775, 2.0).

  Correct values for each bin are: (0.5p for each correct $n$)

  \begin{center}
  \begin{tabular}{rrr}
  & Center of the bin & Number of the galaxies\\
  1 & 0.312500 & 624\\
  2 & 0.537500 & 451\\
  3 & 0.762500 & 354\\
  4 & 0.987500 & 293\\
  5 & 1.212500 & 253\\
  6 & 1.437500 & 210\\
  7 & 1.662500 & 174\\
  8 & 1.887500 & 170
  \end{tabular}
  \end{center}
  % FIGURE NEEDED: histogram of number of galaxies vs distance (MPC),
  % bin size 0.225 MPC (Figure 5, 2022_da_2_sol.pdf page 8)

  In order to get probability density function each value of histogram
  must be divided by it integral on the whole region. The integral can
  be approximated by finding total area of the rectangles:
  \[ I = \sum_{n=1}^{8}H(r_i)\,d = 569.025 \]
  Where $d = 0.225$ is size of a bin. After dividing each data by $I$ we
  get (0.25p for each correct probability density):

  \begin{center}
  \begin{tabular}{rrrr}
  & Center of the bin & Number of the galaxies & Probability density\\
  1 & 0.312500 & 624 & 1.096613\\
  2 & 0.537500 & 451 & 0.792584\\
  3 & 0.762500 & 354 & 0.622117\\
  4 & 0.987500 & 293 & 0.514916\\
  5 & 1.212500 & 253 & 0.444620\\
  6 & 1.437500 & 210 & 0.369052\\
  7 & 1.662500 & 174 & 0.305786\\
  8 & 1.887500 & 170 & 0.298757
  \end{tabular}
  \end{center}
  % FIGURE NEEDED: probability density of center-satellite distances in
  % the same halo (Figure 6, 2022_da_2_sol.pdf page 9)

\item[(f)] Since histogram equation has following form:
  \[ H(r) = \frac{1}{A + B\cdot r} \]
  We should calculate $1/H(r)$ which gives the following values (0.125
  for each correct $1/H$):

  \begin{center}
  \begin{tabular}{rrrr}
  & $r$ & $H(r)$ & $1/H(r)$\\
  1 & 0.312500 & 624 & 0.001603\\
  2 & 0.537500 & 451 & 0.002217\\
  3 & 0.762500 & 354 & 0.002825\\
  4 & 0.987500 & 293 & 0.003413\\
  5 & 1.212500 & 253 & 0.003953\\
  6 & 1.437500 & 210 & 0.004762\\
  7 & 1.662500 & 174 & 0.005747\\
  8 & 1.887500 & 170 & 0.005882
  \end{tabular}
  \end{center}
  % FIGURE NEEDED: linearization plot 1/H(r) vs r for regression
  % (Figure 7, 2022_da_2_sol.pdf page 9)

  Finding the slope and the intercept of this line (by graph or any
  other method) gives following values:
  \[ A = 0.00065977;\qquad B = 0.00285494 \]
\end{description}

\textbf{Solutions -- Part 2}

\begin{description}
\item[(a)] We can $\phi$ values from Part 1.c to distribute them in 60
  degree bins.
  % sic: "We can phi values" as in source (missing "use")
  In bin 1 we have number of galaxies which fall in (0, 60) degree
  range, in bin 2 -- (60, 120) and so on bin 6 -- (300, 360).

  Part 1.c data in degrees:

  \begin{center}
  \begin{tabular}{rrr}
  & Angle & Number of galaxies\\
  1 & 302.496846 & 291\\
  2 & 353.752245 & 170\\
  3 & 303.884926 & 253\\
  4 & 290.179429 & 8\\
  5 & 115.183442 & 72\\
  6 & 294.241133 & 120\\
  7 & 133.773013 & 302\\
  8 & 73.435830 & 52\\
  9 & 203.897219 & 28\\
  10 & 163.073590 & 40\\
  11 & 131.975076 & 72\\
  12 & 43.675004 & 379\\
  13 & 197.372555 & 82\\
  14 & 32.166754 & 62\\
  15 & 133.738158 & 305\\
  16 & 222.698869 & 293
  \end{tabular}
  \end{center}

  We get the following data for the histogram (0.5p for each correct
  $n$):

  \begin{center}
  \begin{tabular}{rrr}
  & Center of the bin & Number of galaxies\\
  1 & 30 & 441\\
  2 & 90 & 124\\
  3 & 150 & 719\\
  4 & 210 & 403\\
  5 & 270 & 128\\
  6 & 330 & 714
  \end{tabular}
  \end{center}

  Uniform value is total number of galaxies divides by 6 which is
  421.5. Square deviation is $241^\circ$ (1p)
  % FIGURE NEEDED: histogram of number of galaxies vs phi angle,
  % bin size 60 deg (Figure 8, 2022_da_2_sol.pdf page 11)

\item[(b)] A galaxy with the orbital radius $r$ around central mass $M$
  has velocity (2p):
  \[ v = \sqrt{\frac{GM}{r}} \]
  This will cause to give the following error in distance along line of
  sight (2p):
  \[ d = \frac{v}{H} = \frac{1}{H}\sqrt{\frac{GM}{r}} \]
  Angle difference for galaxies at 90 and 270 degrees will be (2p):
  \[ \alpha = \arctan\frac{d}{r}
     = \arctan\frac{1}{H}\sqrt{\frac{GM}{r^3}} \]
  % FIGURE NEEDED: geometric drawing of the line-of-sight correction
  % (Figure 9, 2022_da_2_sol.pdf page 12)

\item[(c)] Galaxies originally in bins of range (60, 120) and (240, 300)
  might move to (120, 180) and (300, 360) respectively. So we should
  subtract correction value found in previous task to the galaxies in
  bins 2 and 5 to see if they move to bins 1 and 4 respectively. (0.25p
  for each correct corrected angle)

  \begin{center}
  \begin{tabular}{rrrr}
  & Original angle & Corrected angle & Number of galaxies\\
  1 & 302.494665 & 276.136451 & 291\\
  2 & 353.752069 & 351.714594 & 170\\
  3 & 303.884728 & 300.199512 & 253\\
  4 & 290.178642 & 290.178642 & 8\\
  5 & 115.184577 & 115.184577 & 72\\
  6 & 294.241835 & 294.241835 & 120\\
  7 & 133.776280 & 125.061576 & 302\\
  8 & 73.437747 & 73.437747 & 52\\
  9 & 203.899065 & 203.899065 & 28\\
  10 & 163.073490 & 160.346665 & 40\\
  11 & 131.977481 & 120.548626 & 72\\
  12 & 43.676971 & 43.676971 & 379\\
  13 & 197.372956 & 197.372956 & 82\\
  14 & 32.167948 & 32.167948 & 62\\
  15 & 133.743596 & 94.077612 & 305\\
  16 & 222.694636 & 222.694636 & 293
  \end{tabular}
  \end{center}

  New corrected distribution looks like this (0.3p for each correct
  $n$):

  \begin{center}
  \begin{tabular}{rrr}
  & Center of the bin & Number of galaxies\\
  1 & 30 & 441\\
  2 & 90 & 429\\
  3 & 150 & 414\\
  4 & 210 & 403\\
  5 & 270 & 419\\
  6 & 330 & 423
  \end{tabular}
  \end{center}

  In this case square deviation equals to $12^\circ$ (1p)
  % FIGURE NEEDED: corrected histogram of number of galaxies vs phi
  % angle, bin size 60 deg (Figure 10, 2022_da_2_sol.pdf page 13)
\end{description}

\textbf{Solutions -- Part 3}

\begin{description}
\item[(a)] With simple combinatorics one can get (0.6p for each correct
  $N$):
  \begin{align*}
  N_{CC} &= N_C(N_C - 1) = 249500\\
  N^{*}_{CS} &= N_C N_S = 500000\\
  N_{CS} &= N_C(N_C - 1)N_S = 249500000\\
  N^{*}_{SS} &= N_C N_S(N_S - 1) = 499500000\\
  N_{SS} &= N_C N_S(N_C - 1)N_S = 249500000000
  \end{align*}

\item[(b)] We can calculate all probability densities using
  center-center and center-satellite densities. Non-normalized values
  (with $c_1 = c_2 = c_3 = 1$) are calculated in the following table
  (0.5p for each correct tuple):

  \begin{center}
  \begin{tabular}{rrrrrrr}
  & $r$ & $\rho^{*}_{CS}$ & $\rho_{CC}$ & $\rho^{*}_{SS}$
    & $\rho_{SS}$ & $\rho_{CS}$\\
  1 & 0.312500 & 1.096613 & 0.406141 & 0.375800 & 17.365137 & 0.975747\\
  2 & 0.537500 & 0.792584 & 0.807863 & 0.337652 & 12.992898 & 0.994072\\
  3 & 0.762500 & 0.622117 & 1.340510 & 0.295110 & 8.653487 & 0.988565\\
  4 & 0.987500 & 0.514916 & 1.134011 & 0.261824 & 8.568351 & 0.837480\\
  5 & 1.212500 & 0.444620 & 0.478630 & 0.239696 & 10.729180 & 0.638199\\
  6 & 1.437500 & 0.369052 & 0.143154 & 0.195787 & 10.112507 & 0.458355\\
  7 & 1.662500 & 0.305786 & 0.073165 & 0.155452 & 8.262279 & 0.352612\\
  8 & 1.887500 & 0.298757 & 0.060971 & 0.168470 & 8.998531 & 0.379962
  \end{tabular}
  \end{center}

  For normalization (same way as in Part 1.d) we get $c_1 = 2.18592$;
  $c_2 = 0.79304$; $c_3 = 0.0516$. Multiplying relevant densities by
  this values gives (0.5p for each correct normalized tuple):

  \begin{center}
  \begin{tabular}{rrrrr}
  & $r$ & $\rho^{*}_{SS}$ & $\rho_{SS}$ & $\rho_{CS}$\\
  1 & 0.312500 & 0.822854 & 0.900750 & 0.770962\\
  2 & 0.537500 & 0.739325 & 0.673957 & 0.785441\\
  3 & 0.762500 & 0.646175 & 0.448866 & 0.781089\\
  4 & 0.987500 & 0.573292 & 0.444450 & 0.661713\\
  5 & 1.212500 & 0.524839 & 0.556535 & 0.504257\\
  6 & 1.437500 & 0.428697 & 0.524548 & 0.362158\\
  7 & 1.662500 & 0.340380 & 0.428574 & 0.278607\\
  8 & 1.887500 & 0.368883 & 0.466764 & 0.300217
  \end{tabular}
  \end{center}
  % FIGURE NEEDED: plots of the normalized densities rho_CC, rho_CS,
  % rho_SS (different halos) and rho*_SS (same halo)
  % (Figures 11-14, 2022_da_2_sol.pdf pages 15-16)

\item[(c)] Total probability density is weighted sum of all densities
  (4p):
  \[ \rho \sim \frac{N_{CC}\rho_{CC} + N^{*}_{CS}\rho^{*}_{CS}
     + N_{CS}\rho_{CS} + N^{*}_{SS}\rho^{*}_{SS} + N_{SS}\rho_{SS}}
     {N_{CC} + N^{*}_{CS} + N_{CS} + N^{*}_{SS} + N_{SS}} \]

  Calculated values and normalized densities are given in the following
  table:

  \begin{center}
  \begin{tabular}{rrrr}
  & $r$ & Calculated $\rho$ & Normalized $\rho$\\
  1 & 0.312500 & 0.004212 & 0.843662\\
  2 & 0.537500 & 0.003606 & 0.722433\\
  3 & 0.762500 & 0.002967 & 0.594357\\
  4 & 0.987500 & 0.002693 & 0.539524\\
  5 & 1.212500 & 0.002663 & 0.533361\\
  6 & 1.437500 & 0.002265 & 0.453652\\
  7 & 1.662500 & 0.001813 & 0.363239\\
  8 & 1.887500 & 0.001968 & 0.394216
  \end{tabular}
  \end{center}
  % FIGURE NEEDED: final probability density plot (Figure 15,
  % 2022_da_2_sol.pdf page 16)

  Most similar to this density is plot E on Figure 1 (1p).
\end{description}
